% This is samplepaper.tex, a sample chapter demonstrating the
% LLNCS macro package for Springer Computer Science proceedings;
% Version 2.20 of 2017/10/04
%
\documentclass[runningheads]{llncs}
%
\usepackage{graphicx}
% Used for displaying a sample figure. If possible, figure files should
% be included in EPS format.
%
% If you use the hyperref package, please uncomment the following line
% to display URLs in blue roman font according to Springer's eBook style:
% \renewcommand\UrlFont{\color{blue}\rmfamily}

\begin{document}
%
\title{Contribution Title\thanks{Supported by organization x.}}
%
%\titlerunning{Abbreviated paper title}
% If the paper title is too long for the running head, you can set
% an abbreviated paper title here
%
\author{Kevin Jivani}
%
\authorrunning{F. Author et al.}
% First names are abbreviated in the running head.
% If there are more than two authors, 'et al.' is used.
%
\institute{Concordia University, Montreal, CA}

\maketitle              % typeset the header of the contribution
%
\begin{abstract}
The abstract should briefly summarize the contents of the paper in
15--250 words.

\keywords{First keyword  \and Second keyword \and Another keyword.}
\end{abstract}
%
%
%

\tableofcontents

\section{INTRODUCTION}

In the intricate landscape of project management, the initial steps are often the most critical, determining the course of the entire venture. Three pivotal factors play a paramount role in establishing control over a big project: Project Planning, Project Decomposition, and Effective Communication.

\section{Problem Statement}

\begin{itemize}
  \item How can the decomposition of big projects contribute to establishing control in the initial stages?

The paper discusses the challenges of managing complex design projects and the need for new tools and techniques to aid design managers in making decisions that can reduce the time and cost of a design cycle. As a design manager myself, I can relate to the challenges highlighted in the paper and the need for new tools and techniques to improve the design process.

The paper introduces DeMAID, a tool that was initially released in 1989, and its use in the decision-making process. The tool is designed to aid in the intelligent decomposition of a complex design problem. Since its release, much research has been done in the areas of decomposition, concurrent engineering, and process management, leading to the development of many new tools and techniques.

As a design manager, I am always looking for ways to optimize the design process and reduce the time and cost of a design cycle. The paper highlights the use of graph theory in process-management problems and the network-based approaches, including the PERT and the critical path method. These approaches can be used to reduce the complexity of numerous processes that are dependent on one another.

The paper then goes on to discuss the enhancements made to DeMAID, including the use of a genetic algorithm (GA). The GA is used to optimize the decomposition process by selecting the best decomposition strategy for a given design problem. As a design manager, I am always looking for ways to optimize the design process, and the use of a GA to minimize cost and time in a design cycle is an exciting development.

The paper provides an example of how these enhancements can be applied to improve the design cycle, using an aircraft conceptual design project as an example. The paper explains how the GA was used to optimize the decomposition process, resulting in a significant reduction in the time and cost of the design cycle. As a design manager, I can see the potential of these tools and techniques to revolutionize the design process and improve the efficiency of the design cycle.

In conclusion, the paper highlights the importance of managing complex design projects and the need for new tools and techniques to aid design managers in making decisions that can reduce the time and cost of a design cycle. As a design manager, I can relate to the challenges highlighted in the paper and the need for new tools and techniques to improve the design process. The paper introduces DeMAID and its enhancements, including the use of a genetic algorithm, and provides an example of how these enhancements can be applied to improve the design cycle. The paper provides valuable insights into the use of these tools and techniques and their potential to revolutionize the design process. As a design manager, I am excited about the potential of these tools and techniques to improve the efficiency of the design cycle and reduce the time and cost of a design project.\cite{ref1}

-------------
This research paper proposes a process for defining the necessary structure to reach project objectives and deliver final results. The authors, Julie Stal-Le Cardinal and Franck Marle, argue that building a correct project structure is achievable and gives more guarantee for success, while a bad structure is a guarantee for failure. They suggest that formalizing and managing the project structure definition process is feasible using innovative concepts and tools, such as interactions management, decomposition process, and resource assignment process.

The paper is divided into three parts. Part 1 describes the proposed process, including its inputs, tools, and methods. Part 2 focuses on the inputs of the process, such as the initial situation, objectives, and environment, and presents research proposals on environment management. Part 3 presents research proposals on tools and methods, specifically on the topics of scope and activity definition and resource assignment. The authors provide concrete examples from their work with VALLOUREC and PSA PEUGEOT CITROEN.

The methodology used in this paper is based on a literature review of existing research on project management, as well as the authors' own experience and research. The authors also draw on their work with VALLOUREC and PSA PEUGEOT CITROEN to provide practical examples of the proposed process and tools. Overall, the paper provides a comprehensive framework for managing projects and designing project structures that can help ensure success. \cite{ref2}

  
  \item \textbf{What role does leadership play in initiating and planning large programs for effective control?}
  \item How does the planning process contribute to program control, especially in the context of managing interproject dependencies?
  \item What communication strategies and infrastructure are essential for managing interactions and maintaining control in a large program?
\end{itemize}

\subsection{Project Planning}

Project management in the context of large-scale endeavors is a complex and challenging undertaking that requires a strategic approach from the very beginning. The initial steps in establishing control over a massive project involve meticulous Project Planning. This phase marks the inception of the project management life cycle, where the business problem or opportunity is identified, and a comprehensive solution is outlined. A crucial component of this stage is the creation of a business case, which serves as a detailed roadmap, encompassing problem statements, alternative solutions, cost-benefit analysis, and implementation plans. The approval of the business case triggers the transition to the planning phase, emphasizing the significance of having a clear understanding of the project's objectives. Without a well-defined plan, projects risk deviating from their intended course, leading to misunderstandings and potential failures. Establishing clear objectives is akin to a compass guiding the project towards success, ensuring that everyone involved shares a unified vision.

\subsection{Project Decomposition}

Once the project plan is in place, the next pivotal factor is Project Decomposition. This involves breaking down the colossal project into smaller, more manageable components. Whether it's subdividing tasks or disassembling a large product into smaller parts, decomposition project management enhances manageability and facilitates improved communication and coordination among team members. Unlike traditional linear project management models, which may become unwieldy as complexity grows, decomposition takes a flexible approach. It breaks the project into smaller, independent pieces, allowing for greater adaptability and a comprehensive understanding of the project's entirety. The process involves establishing the big picture, defining major objectives, breaking them into manageable chunks, and creating a work breakdown structure. This systematic breakdown not only eases the workload but also enables better error identification and rectification as each component is assessed independently.

\subsection{Effective Communication}

In the realm of project management, where intricate tasks and numerous responsibilities coalesce, Effective Communication emerges as the linchpin for success. Communication is not merely a peripheral aspect but an integral component that ensures everyone involved comprehends the project's goals and responsibilities. In the vast landscape of project management, effective communication skills wield immense power. A project manager, regardless of the project's scale, must articulate their vision clearly, ensuring it resonates with the entire team. In larger companies, where projects span multiple departments, communication becomes even more critical. The project manager must not only convey their vision for a specific segment but also collaborate with other managers to ensure the seamless completion of the overarching project. The significance of communication lies not only in transmitting vital information but also in cultivating a shared vision among team members. It involves clearly defining project goals, ensuring understanding among team members, and outlining performance expectations. Active communication methods, such as in-person meetings, video conferencing, and webinars, become invaluable tools for fostering a collaborative and informed project environment. Tailoring communication methods to the project's specific needs ensures that everyone involved receives information in a manner that suits them best.

% In essence, project management is a multifaceted discipline that demands a holistic approach. Beyond effective communication, successful project managers possess a diverse set of skills. These encompass the ability to connect tasks to the bigger picture, prioritize goals, maintain an agile mindset, and react adeptly to change. While basic communication skills are essential, they are only one facet of a project manager's repertoire. The value of project management to companies is evident in its capacity to enhance organizational efficiency, save costs, align projects with company needs, and provide a competitive edge. Aspiring project managers can develop these skills through comprehensive programs such as the online Master of Communication Management (MCM), which not only hones communication abilities but also instills a strategic, holistic approach to project management. In the dynamic landscape of large projects, the combination of effective planning, meticulous decomposition, and clear communication emerges as the trinity that propels success.






\section{Motivation}
Understanding how to control a project from the outset is paramount, especially when dealing with substantial endeavors. A massive project brings with it a multitude of complexities and intricacies that, if not managed effectively, can lead to significant challenges. Project control involves the systematic oversight and coordination of various elements, ensuring that the project stays on course and aligns with its objectives. In the early stages of a large project, establishing control mechanisms is like putting a solid foundation in place; it provides stability and sets the tone for the entire project lifecycle.

One primary reason for emphasizing control in the initial phases is the prevention of scope creep. In a large project, the scope, or the defined boundaries of what the project aims to achieve, can easily expand beyond the original plan. Without proper control mechanisms, additional features or requirements might be introduced, leading to increased costs, delayed timelines, and potential resource strain. Control allows project managers to delineate and communicate the project's scope clearly, creating a framework that helps prevent unnecessary deviations.

Furthermore, control is essential for resource allocation and utilization. In substantial projects, various resources, including human, financial, and technological, are at play. Efficient control ensures that these resources are allocated optimally, avoiding overloading specific areas while neglecting others. By understanding and implementing control early on, project managers can identify potential resource bottlenecks and address them proactively, promoting a balanced and effective distribution of resources throughout the project lifecycle.

Risk management is another critical aspect tied to project control. Large projects inherently come with an increased level of uncertainty and potential risks. Without a robust control framework, these risks may escalate, jeopardizing the project's success. Control mechanisms empower project managers to identify, assess, and mitigate risks systematically. By acknowledging and addressing risks early in the project, managers can make informed decisions, implement preventive measures, and ensure that the project stays resilient in the face of uncertainties.

Additionally, effective control contributes to stakeholder satisfaction. In extensive projects, numerous stakeholders with varying interests and expectations are involved. Clear communication, transparency, and regular monitoring foster trust among stakeholders. When project managers have control mechanisms in place, they can provide stakeholders with accurate updates on project progress, potential challenges, and mitigation strategies. This transparency builds confidence and collaboration, aligning the project team and stakeholders toward shared objectives.

% \section{Background Material}





% \subsection{A Subsection Sample}
% Please note that the first paragraph of a section or subsection is
% not indented. The first paragraph that follows a table, figure,
% equation etc. does not need an indent, either.

% Subsequent paragraphs, however, are indented.

% \subsubsection{Sample Heading (Third Level)} Only two levels of
% headings should be numbered. Lower level headings remain unnumbered;
% they are formatted as run-in headings.

% \paragraph{Sample Heading (Fourth Level)}
% The contribution should contain no more than four levels of
% headings. Table~\ref{tab1} gives a summary of all heading levels.

% \begin{table}
% \caption{Table captions should be placed above the
% tables.}\label{tab1}
% \begin{tabular}{|l|l|l|}
% \hline
% Heading level &  Example & Font size and style\\
% \hline
% Title (centered) &  {\Large\bfseries Lecture Notes} & 14 point, bold\\
% 1st-level heading &  {\large\bfseries 1 Introduction} & 12 point, bold\\
% 2nd-level heading & {\bfseries 2.1 Printing Area} & 10 point, bold\\
% 3rd-level heading & {\bfseries Run-in Heading in Bold.} Text follows & 10 point, bold\\
% 4th-level heading & {\itshape Lowest Level Heading.} Text follows & 10 point, italic\\
% \hline
% \end{tabular}
% \end{table}


% \noindent Displayed equations are centered and set on a separate
% line.
% \begin{equation}
% x + y = z
% \end{equation}
% Please try to avoid rasterized images for line-art diagrams and
% schemas. Whenever possible, use vector graphics instead (see
% Fig.~\ref{fig1}).

% \begin{figure}
% \includegraphics[width=\textwidth]{fig1.eps}
% \caption{A figure caption is always placed below the illustration.
% Please note that short captions are centered, while long ones are
% justified by the macro package automatically.} \label{fig1}
% \end{figure}

% \begin{theorem}
% This is a sample theorem. The run-in heading is set in bold, while
% the following text appears in italics. Definitions, lemmas,
% propositions, and corollaries are styled the same way.
% \end{theorem}
% %
% % the environments 'definition', 'lemma', 'proposition', 'corollary',
% % 'remark', and 'example' are defined in the LLNCS documentclass as well.
% %
% \begin{proof}
% Proofs, examples, and remarks have the initial word in italics,
% while the following text appears in normal font.
% \end{proof}
% For citations of references, we prefer the use of square brackets
% and consecutive numbers. Citations using labels or the author/year
% convention are also acceptable. The following bibliography provides
% a sample reference list with entries for journal
% articles~\cite{ref_article1}, an LNCS chapter~\cite{ref_lncs1}, a
% book~\cite{ref_book1}, proceedings without editors~\cite{ref_proc1},
% and a homepage~\cite{ref_url1}. Multiple citations are grouped
% \cite{ref_article1,ref_lncs1,ref_book1},
% \cite{ref_article1,ref_book1,ref_proc1,ref_url1}.
% %
% % ---- Bibliography ----
% %
% % BibTeX users should specify bibliography style 'splncs04'.
% % References will then be sorted and formatted in the correct style.
% %
% % \bibliographystyle{splncs04}
% % \bibliography{mybibliography}
%
\begin{thebibliography}{8}
\bibitem{ref_article1}
Author, F.: Article title. Journal \textbf{2}(5), 99--110 (2016)

\bibitem{ref1}
Tools and Techniques for Decomposing and Managing Complex Design Projects
James L. Rogers
Journal of Aircraft 1999 36:1, 266-274  \doi{10.2514/2.2434}

\bibitem{ref2}
Julie Stal-Le Cardinal, Franck Marle,
Project: The just necessary structure to reach your goals,
International Journal of Project Management,
Volume 24, Issue 3,
2006,
Pages 226-233,
ISSN 0263-7863,
https://doi.org/10.1016/j.ijproman.2005.10.002.

\bibitem{ref_proc1}
Author, A.-B.: Contribution title. In: 9th International Proceedings
on Proceedings, pp. 1--2. Publisher, Location (2010)

\bibitem{ref_url1}
LNCS Homepage, \url{http://www.springer.com/lncs}. Last accessed 4
Oct 2017
\end{thebibliography}
\end{document}
